\section{Conclusions and Future Work}
\label{sec:conclusions}

\subsection{Conclusions}
\label{subsec:conc-conclusions}

This project set out to build, analyse and evaluate a cooperative multi-agent
system for a deceptively simple task, finding a hidden target on a grid, made
genuinely hard by three coupled sources of uncertainty: partial observability,
limited communication, and, as its defining extension, the possibility that any
drone permanently breaks at any moment with no central authority to announce the
loss.

\paragraph{What was built, and what the multi-agent reading revealed.}
We implemented from scratch a cooperative grid environment with sequential
execution, accumulated per-drone knowledge maps, range-limited fusion, and a
decentralized fault model in which crash knowledge propagates only through a
passive distress beacon and ordinary communication; a single Dueling Double DQN
with parameter sharing drives every drone, and per-episode domain randomization
with a compact context vector makes that one policy a generalist across team
size, sensing, communication, clutter and fault rate. Read through the vocabulary
of the field (Section~\ref{sec:mas}), the system's most interesting properties
turn out to be precisely the ones no line of code states directly. Coordination
is achieved \emph{without} protocols, performatives or a coordinator: dispersion
from a shared spawn corner is driven by environment-mediated, stigmergic
repulsion from swept regions; role differentiation among identical drones is
broken open by a single identity scalar; and the ``knowledge'' each drone acts on
is a first-order epistemic belief, sound but possibly stale, updated by
observation and by testimony from teammates. Graceful degradation, in particular,
is a system-level property that no individual drone implements, resting on
redundancy and epistemic decentralization rather than on any drone knowing the
true state of the team.

\paragraph{What the experiments and the ethical analysis showed.}
The evaluation (Section~\ref{sec:experiments}) substantiated the central claims
and bounded them honestly. A single policy generalized across every randomization
axis, including \emph{beyond} its trained ranges on team size, vision and
communication, driving teams of six drones it never trained on at over $90\%$
success. Most importantly, degradation under faults is \emph{graceful}: across
the entire trained fault range success fell by under six points while the team
shed, on average, more than half a drone, and even at a fault rate more than
triple the trained maximum the system still completed two missions in three, with
no cliff. The sweeps also separated two competences the success rate conflates:
vision governs both whether and how well the team searches, while communication
governs only how well, though it governs it substantially, map fusion buying
close to a fifth in path quality, mission time and redundancy at no cost beyond
proximity. The limits were equally clear: obstacle density is the one axis whose
extrapolation collapses; and two of the three emergence signatures, together with
the belief-lag question, remain defined but uninstrumented.
Section~\ref{sec:ethics} then argued that the capability is irreducibly
dual-use, the fault tolerance that makes a search-and-rescue swarm robust being
what would make an attack swarm resilient; mapped the civilian system onto the
GDPR and the EU AI Act; confronted the military carve-out that lets the war
scenario escape the Act entirely; and concluded that \emph{this specific system},
opaque, statistically rather than formally validated, acting on stale beliefs and
with no single decision locus, would be unacceptable as an autonomous weapon, on
exactly the grounds its own evaluation exposes. The takeaway is not that the
technology should not be built, but that its governance cannot be an afterthought
to its engineering.

\subsection{Future work}
\label{subsec:conc-future}

Three measurements were defined in this report but not carried out, each for a
mundane reason. The belief-lag question, that is how communication range governs
the interval between a fault and a survivor's discovery of it, needs a joint
fault$\times$comm sweep with per-step trajectory logging, and the same logging
would deliver the two emergence signatures (inter-drone distance over time,
post-fault re-covering latency) that Section~\ref{subsec:exp-qualitative} left
qualitative; two further ablations, removing the believed-alive count from the
context vector and switching the shared terminal reward for an individual one,
require retraining a policy each. None of these changes the architecture: they
are instrumentation and compute, and they would sharpen the claims already made.

Three directions reach further. On coordination, the finding that communication
buys efficiency but not find rate here is partly a property of the task and
partly of the
design: \emph{learned} messages would let the system decide \emph{what} to share,
and a task that genuinely rewarded information sharing (a target visible only
from afar, heterogeneous sensors, partial views that must be combined) would
exercise that channel far harder than target-in-a-maze search does, while
value-decomposition methods (VDN/QMIX-style) would replace the shared-terminal
credit heuristic with a learned factorization of team value. On representation,
the hand-crafted knowledge maps that let a memoryless policy act under partial
observability could give way to a learned recurrent or attention-based memory
over the observation history. And on the world itself, every limitation of
Section~\ref{subsec:exp-discussion} is a research direction: moving or evasive
targets, correlated or passage-blocking faults, and, most pointedly given the
ethical analysis, an \emph{adversary} choosing when and whom to disable and how
to jam communication, turning the benign i.i.d.\ fault process into a strategic
one and the cooperative game into a mixed one. Closing the sim-to-real gap while
preserving the fault tolerance that is the point of the exercise is the direction
in which the work would have to move to matter outside the laboratory.
